\section{From random charts to uniform height control}
\label{sec:uniformity}

The preceding sections identify affine profiles and compare nearby
sheets.  We now remove the random chart center.  There are two distinct
points to check: the affine label must be independent of scale and
basepoint, and the chart offset must disappear after division by the
scale.  Throughout this section the Lipschitz, quasi-isometry and
coarse-inverse bounds for the pairs $(F,\bar F)$ are fixed.  All constants
are uniform over that family.

Choose $\delta>0$ small enough for Proposition~\ref{prop:neighbor-offsets},
and use $\mathcal C$, $\mathcal U_\delta$, and
$f_{r,x,n}$ from Definition~\ref{def:chart-profiles}.  Thus
\[
 f_{r,x,n}(z)=
 \frac{\pi F(xn\Delta_r z)-\pi F(x)}r,
 \qquad z\in D_\infty.
\]
Choose a compact identity neighborhood $Z\subset\mathcal U_\delta$.
In particular, $\pi Z$ contains a neighborhood of zero in $W$.
We can shrink $Z$ a fixed number of times below.

\begin{lemma}[Suffix comparison and concentration]
\label{lem:suffix-concentration}
Let $a$ be uniform on $B_r(\mathcal C)$.  Along every sequence
$r\to\infty$, with arbitrary maps $F$, basepoints $x$, and suffixes
\[
 k_0\in B_r((1+\delta)\mathcal C,e^{\eta_r r}),
 \qquad \eta_r\longrightarrow0,
\]
one has
\begin{equation}
\label{eq:suffix-profile-comparison}
 \|f_{r,x,ak_0}-f_{r,x,a}\|_Z\longrightarrow0
 \quad\text{in probability}.
\end{equation}
The assertion also holds for suffixes sampled independently of $a$
whose supports satisfy the same bound.

There are deterministic $A_{F,x,r}\in\mathcal A$ and $s_{F,x,r}\in W$
such that
\begin{equation}
\label{eq:ordinary-chart-concentration}
 \sup_{F,x}\mathbb E\left\|
 f_{r,x,a}-A_{F,x,r}\pi-\frac{s_{F,x,r}}r
 \right\|_Z\longrightarrow0,
 \qquad |s_{F,x,r}|\le C r.
\end{equation}
The same approximation holds in probability with $a$ replaced by
$ak_0$ in the sequential setting of
\eqref{eq:suffix-profile-comparison}.
\end{lemma}

\begin{proof}
By Lemma~\ref{lem:product-schemes}, an allowed suffix is a product
of a bounded number of generator exponentials.  Their amplitudes obey
the bounds of Proposition~\ref{prop:neighbor-offsets}; all partial
products remain allowed suffixes by Lemma~\ref{lem:box-calculus}.
The joint law of each successive pair of profiles is relatively
compact by Proposition~\ref{prop:affine-profiles}.  Every limit is
supported on identical pairs by Proposition~\ref{prop:neighbor-offsets}.
The distance between the two profiles on $Z$ therefore tends to zero
in probability.  Summing over the bounded number of steps proves
\eqref{eq:suffix-profile-comparison}.  If independent random suffixes
violated the assertion, conditioning would select deterministic
suffixes with failure probability bounded below.  This would contradict
the deterministic statement.

Write $\mu_r$ for uniform Haar probability on $B_r(\mathcal C)$.
The normalized product bounds give a fixed $L\ge1$ such that
\[
 a^{-1}B_r(\mathcal C)\subset B_r(\mathcal C,L)
 \qquad(a\in B_r(\mathcal C)).
\]
Let $k$ be independent and uniform on $B_r(\mathcal C,L)$.
For each fixed $a$, the conditional law of $ak$ dominates $c\mu_r$
for a constant $c>0$ independent of $r$ and $a$.  Indeed, its Haar
density on $aB_r(\mathcal C,L)$ is the reciprocal of the volume of
that box, and the ratio of this volume to $\vol B_r(\mathcal C)$
is bounded.  Consequently, for independent $a,b$ with law $\mu_r$,
\[
 \Pr\{\|f_{r,x,a}-f_{r,x,b}\|_Z>\epsilon\}
 \le c^{-1}
 \Pr\{\|f_{r,x,a}-f_{r,x,ak}\|_Z>\epsilon\}
 \longrightarrow0.
\]

The ordinary profile laws are relatively compact and uniformly bounded
on $Z$.  If $\nu$ is any subsequential limit, the last display says that
$\nu\times\nu$ is supported on the diagonal.  Thus $\nu$ is a point
mass.  Proposition~\ref{prop:affine-profiles} says that this point has
the form $A\pi+b$, with $A\in\mathcal A$.

Set
\[
 s_{F,x,r}=r\,\mathbb E f_{r,x,a}(1),
\]
and choose $A_{F,x,r}\in\mathcal A$ to minimize the expected supremum
norm in \eqref{eq:ordinary-chart-concentration}.  These quantities
are well defined because the profiles are bounded.  The point-mass
conclusion proves convergence to zero along every sequence of maps
and basepoints; otherwise a violating sequence would have a
point-mass subsequential limit for which the minimum is zero.
This proves the stated supremum over $F,x$.
Finally $|a|_G\le Cr$ by Lemma~\ref{lem:log-distortion}, so the
Lipschitz bound for $F$ gives $|s_{F,x,r}|\le Cr$.
Equation~\eqref{eq:suffix-profile-comparison} transfers the deterministic
approximation to the allowed suffixes.
\end{proof}

\begin{remark}
\label{rem:suffix-envelope}
The suffix assertion is sequential.  Equivalently, one can prescribe
a vanishing envelope $\eta(r)$ for any fixed finite collection of
constructions and then take the corresponding uniform bounds.  It
does not mean a supremum at a fixed $r$ over every sequence that might
eventually satisfy $\eta_r\to0$.  The supremum in
\eqref{eq:ordinary-chart-concentration} concerns ordinary charts and
has no such suffix parameter.
\end{remark}

\begin{lemma}[Stabilization across scales]
\label{lem:scale-stabilization}
There is $R_0$ independent of $F,x$ such that $A_{F,x,r}$ is independent
of $r\ge R_0$.  Moreover
\[
 \sup_{F,x}|s_{F,x,r}|/r\longrightarrow0.
\]
\end{lemma}

\begin{proof}
Put $q=1+\delta/2$.  If $r\le r'\le qr$, let $a$ be uniform on
$B_r(\mathcal C)$ and let $k$ be independent and uniform on
$B_{r'}(\mathcal C,L)$, for a sufficiently large fixed $L$.
For every such $a$,
\[
 a^{-1}B_{r'}(\mathcal C)\subset B_{r'}(\mathcal C,L).
\]
The density argument in Lemma~\ref{lem:suffix-concentration} shows
that the law of $ak$ dominates a fixed multiple of uniform measure
on $B_{r'}(\mathcal C)$.  This suffix fits the strict expansion at
scale $r$, since $r'\le(1+\delta/2)r$.

For the same $N$-coordinate $n$, compare the two deterministic
approximations at
\[
 z'=\Delta_{r'}^{-1}\Delta_rz,
 \qquad r'\pi z'=r\pi z.
\]
On a fixed smaller $Z$, $z'$ stays in the original neighborhood.
The errors of both approximations are $o(r)$ in probability,
uniformly over $r'\in[r,qr]$, $F$, and $x$.  To justify this last
uniformity, a violating sequence would be an allowed sequence in
Lemma~\ref{lem:suffix-concentration}, with the fixed envelope coming
from $L$; the density comparison constants are fixed.
More precisely, the lower domination of the law of $ak$ transfers
its scale-$r$ good event to uniform measure on $B_{r'}(\mathcal C)$.
Intersect there with the ordinary scale-$r'$ good event.
Therefore
\begin{equation}
\label{eq:scale-affine-comparison}
 \sup_{z\in Z}
 \left|s_{F,x,r}-s_{F,x,r'}
       +r(A_{F,x,r}-A_{F,x,r'})\pi z\right|=o(r).
\end{equation}
The expression is deterministic: an event of positive probability
on which both approximations are accurate proves its bound.

Subtract \eqref{eq:scale-affine-comparison} at $1$ from its values
at a fixed finite set of points whose heights span $W$.
The offsets cancel.  Distinct members of the finite set $\mathcal A$
are separated on these heights, so the labels agree for all
$r\ge R_0$, with one $R_0$.  Since any larger scale can be reached
by successive ratios in $[1,q]$, the label stabilizes.
Evaluating at $1$ also gives
\[
 |s_{F,x,r'}-s_{F,x,r}|\le\varepsilon(r)r
 \quad(r\le r'\le qr),\qquad \varepsilon(r)\longrightarrow0,
\]
uniformly in the other variables.

For completeness, fix $\epsilon>0$ and choose a fixed $R\ge R_0$
such that $\varepsilon(t)\le\epsilon$ for $t\ge R$.
Telescope along scales $R,qR,q^2R,\ldots$, ending with a ratio
at most $q$.  The uniform bound at the initial scale gives
\[
 |s_{F,x,r}|
 \le CR+\epsilon\sum_{j:q^jR<r}q^jR
 \le CR+\frac{q}{q-1}\epsilon r.
\]
First divide by $r$ and let $r\to\infty$, then let $\epsilon\to0$.
This proves the uniform sublinear bound.
\end{proof}

\begin{proposition}[Uniform sublinear height estimate]
\label{prop:uniform-sublinear}
For each map $F$ there is a single $A_F\in\mathcal A$ and there is
a function $\epsilon(R)\to0$, common to the entire fixed family,
such that
\begin{equation}
\label{eq:uniform-sublinear}
 \left|\pi F(xg)-\pi F(x)-A_F\pi g\right|
 \le\epsilon(R)R
 \qquad(F,\ x\in G,\ |g|_G\le R).
\end{equation}
The label of the coarse inverse is $A_{\bar F}=A_F^{-1}$.
\end{proposition}

\begin{proof}
It remains to compare the stabilized labels and center values at
different basepoints.  Fix a sufficiently small $\eta>0$ relative
to $\delta$.  Suppose $|x^{-1}y|_G\le\eta r$ and choose a lift
\[
 x^{-1}y=ud,\qquad
 |d|_D\le C(1+\eta r),\qquad
 \|\log u\|\le e^{C(1+\eta r)}
\]
using Lemma~\ref{lem:path-packing}.  Set $r_0=(1-\delta/2)r$ and,
for $a$ uniform on $B_r(\mathcal C)$, put
\[
 a_y=d^{-1}u^{-1}ad.
\]
Its law has constant Haar density $1/\vol B_r(\mathcal C)$ on
its support, since conjugation by $d$ has determinant one.

Let $a_0$ be uniform on $B_{r_0}(\mathcal C)$ and let $k$ be
independent and uniform on
$d^{-1}B_r(\mathcal C,L)d$.  We claim that the law of $a_y$
is dominated by a fixed multiple of the law of $a_0k$.
Indeed, for every possible $a,a_0$,
\[
 \Ad_d(a_0^{-1}a_y)=(\Ad_da_0)^{-1}u^{-1}a
 \in B_r(\mathcal C,L).
\]
The inclusion follows from Lemma~\ref{lem:box-calculus}:
the smaller box retains a strict exponent gap after conjugation
by $d$, and $u$ fits strictly within the $r$-widths if $\eta$
is small enough.  Polynomial conjugation factors are absorbed by
these gaps.  The Haar volume of the box for $k$ is a fixed
multiple of $\vol B_r(\mathcal C)$.  Conditional density comparison,
followed by integration in $a_0$, proves the claim.

All these $k$ lie in the allowed expansion at scale $r_0$.
For this use
$(1+\delta)(1-\delta/2)>1$ and choose $\eta$ smaller if necessary;
the remaining strict gap again absorbs the polynomial factors.
Thus Lemma~\ref{lem:suffix-concentration} applies with a common
envelope to this collection of comparisons.

We have the exact identity
\[
 xa\Delta_rz=ya_y d^{-1}\Delta_rz.
\]
For $z$ in a fixed smaller neighborhood, the point
$z'=\Delta_{r_0}^{-1}(d^{-1}\Delta_rz)$ lies in $Z$.
This follows from Lemma~\ref{lem:d-scaling} and the small bound
on $|d|_D/r$; its height satisfies exactly
\[
 r_0\pi z'=r\pi z-\pi d.
\]
The two deterministic chart approximations, transferred through
the fixed density domination, can now be compared.  Their offsets
are $o(r)$ by Lemma~\ref{lem:scale-stabilization}.  Subtracting the
comparison at two values of $z$ cancels both center values and
$\pi d$.  Testing a fixed spanning set of heights forces equality
of the stabilized labels at $x,y$, uniformly for all sufficiently
large $r$.  Evaluation at $z=1$ then gives
\[
 \pi F(y)-\pi F(x)-A_{F,y,r_0}\pi(x^{-1}y)=o(r),
 \qquad |x^{-1}y|_G\le\eta r,
\]
uniformly in $F,x,y$.

Any two basepoints can be compared at a sufficiently large scale;
therefore the stabilized label is a single $A_F$.
To obtain \eqref{eq:uniform-sublinear} for $|g|\le R$, take
$y=xg$ and $r=R/\eta$.  The preceding uniform error then provides
one modulus $\epsilon(R)\to0$ for all maps and basepoints.

Finally apply \eqref{eq:uniform-sublinear} to both maps in a
coarse-inverse pair.  For $g_t=\exp_D(tH)$, the displacement
$F(g_t)$ from $F(1)$ is $O(t)$.  Bounded-error composition gives
\[
 t\pi H
 =A_{\bar F}A_F(t\pi H)+o(t).
\]
Every $w\in W$ is $\pi H$ for some $H\in\mathfrak d$.
Division by $t$ proves $A_{\bar F}A_F=\id$.
\end{proof}
